\begin{table}[htbp]
\centering
\small
\caption{Observed against null spatial concentration of the lowest accessibility decile.}
\label{tab:t15_null}
\begin{tabular}{llrrrlll}
\toprule
City & Flagged enterprises & Zones, observed & Zones, null & In zones, observed (\%) & In zones, null (\%) & $z$ & $p$ \\
\midrule
Kochi & 2,777 & 5 & 7.400000 & 84.800000 & 62.9 [54.7–68.9] & +4.67 & 0.000 \\
Varanasi & 929 & 5 & 1.200000 & 60.900000 & 24.7 [0.0–53.8] & +1.52 & 0.030 \\
Kolkata & 2,489 & 14 & 3.600000 & 92.000000 & 91.2 [56.9–100.0] & +0.05 & 0.790 \\
Patna & 1,290 & 8 & 2.400000 & 67.200000 & 90.1 [58.0–99.4] & -1.78 & 0.890 \\
Indore & 4,892 & 18 & 2.600000 & 67.700000 & 86.9 [80.9–92.9] & -4.36 & 1.000 \\
\bottomrule
\end{tabular}
\begin{tablenotes}\footnotesize
\item Null is 100 random subsets of equal size drawn from the same enterprise population; the bracketed range is the 5th to 95th percentile. MFAI is a smooth function of distance to fixed terminals, so low-accessibility regions are contiguous by construction and flagged enterprises cluster whatever the threshold; excess concentration over the null is the evidence. In Patna and Indore flagged enterprises are significantly \emph{less} concentrated than an arbitrary equally sized group, because they are peripheral and the periphery is sparse. $p$ is one-sided.
\end{tablenotes}
\end{table}
